%---------------------------------------------------------------
% Eigene Befehle, Makros, Abkürzungen usw.
% Einige erfordern besondere Packages!
%
% Im Allgemeinen keine Neudefinition bekannter Befehle.
% Ausnahmen:
% \t Tabelle
% \u steile Einheit mit Abstand
% \dh Abkürzung für 'das heißt'
%---------------------------------------------------------------

%---------------------------------------------------------------
% Allgemeines
%---------------------------------------------------------------
\newcommand{\nc}{\newcommand}
\nc{\rnc}{\renewcommand}
%
%---------------------------------------------------------------
% Textformatierungen
%---------------------------------------------------------------
\nc{\cl}{\centerline}
\nc{\x}[1]{\mbox{#1}}                   % mbox zur Vermeidung von Trennungen
%\nc{\figfont}{\normalsize \rm}
%\nc{\bvb}{\begin{verbatim}}
%\nc{\evb}{\end{verbatim}}
%
%---------------------------------------------------------------
% Bildunterschriften: 'a) Text' usw. nicht trennen
%---------------------------------------------------------------
\nc{\xa}{: a)~}
\nc{\xb}{; b)~}
\nc{\xc}{; c)~}
\nc{\xd}{; d)~}
\nc{\xe}{; e)~}
\nc{\xf}{; f)~}
%
%---------------------------------------------------------------
% Abstände
%---------------------------------------------------------------
\nc{\hs}[1]{\hspace*{#1}}
\nc{\vs}[1]{\vspace*{#1}}
%
%---------------------------------------------------------------
% Umschaltung Sprache für verschiedene Trennverzeichnisse
% (interessant z.B. für Lit.verzeichnisse, in der Datenbasis
% sollte \de bzw. \en enthalten sein).
%---------------------------------------------------------------
\nc{\en}{\selectlanguage{USenglish}}
\nc{\de}{\selectlanguage{ngerman}}
%
%---------------------------------------------------------------
% Abkürzungen für Literaturdatenbank
%---------------------------------------------------------------
\nc{\fb}[1]{. {\de \rm Fortschrittberichte VDI, \x{Reihe 11}: Schwingungstechnik, \x{Nr. #1}}}
\nc{\te}[1]{, \en Int.~Gas Turbine~\&~Aeroeng.~Congress~\&~Exh., #1}
\nc{\jsv}{\en J.~of Sound and Vibration}
\nc{\vdiv}{\de VDI-Verlag}
\nc{\ddorf}{\de D{\"u}sseldorf}
%
%---------------------------------------------------------------
% Literatur zitieren
%---------------------------------------------------------------
\nc{\etal}{et~al.}
\nc{\z}[2]{\x{\sc #1}~{\x{\cite{#2}}}}
\nc{\zeigene}[2]{\x{\sc #1}~{\x{\citeeigene{#2}}}}
\nc{\zandere}[2]{\x{\sc #1}~{\x{\citeandere{#2}}}}
\nc{\nz}[2]{\x{\sc #1}~{\x{\nocite{#1}}}}
\nc{\nzeigene}[1]{{\x{\nociteeigene{#1}}}} 
\nc{\nzandere}[1]{{\x{\nociteandere{#1}}}}
\nc{\zo}[1]{\x{\cite{#1}}}
\nc{\ztr}[2]{{\sc #1}~{\x{\cite{#2}}}}
%
%---------------------------------------------------------------
% Abkürzungen
%---------------------------------------------------------------
\nc{\zb}{z.\,B.\xspace}
\nc{\zt}{z.\,T.\xspace}
\nc{\ua}{u.\,a.\xspace}
\nc{\uae}{u.\,ä.\xspace}
\nc{\oa}{o.\,a.\xspace}
\nc{\oae}{o.\,ä.\xspace}
\nc{\bzw}{bzw.\xspace}
\nc{\uu}{u.\,U.\xspace}
\rnc{\dh}{d.\,h.\xspace}
\nc{\ie}{i.\,e.\xspace}
\nc{\eg}{e\,g.\xspace}
%
%---------------------------------------------------------------
% Programmnamen und Eigennamen
%---------------------------------------------------------------
% Problem mit \sf unter Verwendung von Type1 Fonts und pdflatex.
% Nach updmap mit der updmap.cfg von Martin steht die sfss1200.pfb,
% die sonst mit der alten Version von updmap.cfg verwendet wurde,
% nicht mehr zur Verfügung. Abhilfe: in updmap.cfg den Mapfile der cm-super-type1-
% schriften mit aufnehmen, dann wird sfss1200.pfb sauber eingebunden. Oder aber
% nicht \sf verwenden, sondern mit \newfont{\ssfont}{cmss10} auf bekannten Type1-
% Font aus der bsr.map zurückgreifen. Problem: Skalierung bei eigener Definition.
%
% Erfordert xspace.sty zur Vermeidung von Leerzeichenfehlern nach eigenen Definitionen
%
\nc{\prog}[1]{{\sf{#1}}\xspace}
%
\nc{\cadax}{\prog{\x{CADA-X}}}
\nc{\matlab}{\prog{Matlab}}
\nc{\simulink}{\prog{Simulink}}
\nc{\datur}{\prog{\x{Datur}}}
\nc{\dates}{\prog{\x{Dates}}}
\nc{\danko}{\prog{\x{Danko}}}
\nc{\datar}{\prog{\x{Datar}}}
\nc{\dankodatar}{\prog{\x{Danko/Datar}}}
\nc{\ansys}{\prog{\x{Ansys}}}
\nc{\hades}{\prog{\x{Hades}}}
\nc{\sco}{\prog{\x{SC03}}}
\nc{\sj}{\prog{\x{SJ18}}}
\nc{\conv}{\prog{CONVERSION}}
%\nc{\sol}{\prog{SOLUTION}}
\nc{\pre}{\prog{PREPROCESSING}}
\nc{\post}{\prog{POSTPROCESSING}}
\nc{\opt}{\prog{OPTIMIZATION}}
\nc{\unix}{\prog{\x{UNIX}}}
\nc{\windows}{\prog{WINDOWS}}
\nc{\ghostview}{\prog{GhostView}}
\nc{\ghostscript}{\prog{GhostScript}}
\nc{\dvips}{\prog{Dvips}}
\nc{\miktex}{\prog{Miktex}}
\nc{\bibtex}{\prog{Bibtex}}
\nc{\powerpoint}{\prog{PowerPoint}}
\nc{\fortran}{\prog{Fortran}}
\nc{\winedt}{\prog{WINEDT}}
\nc{\coreldraw}{\prog{CorelDraw}}
\nc{\prettyplot}{\prog{PrettyPlot}}
\nc{\adistiller}{\prog{AcrobatDistiller}}
\nc{\areader}{\prog{AcrobatReader}}
\nc{\distiller}{\prog{Distiller}}
\nc{\java}{\prog{JAVA}}
\nc{\ascii}{\prog{ASCII}}
%
%---------------------------------------------------------------
% Tex-Specials
%---------------------------------------------------------------
\nc{\tex}{\TeX\xspace}
\nc{\latex}{\LaTeX\xspace}
\nc{\pdflatex}{Pdf\LaTeX\xspace}
%
%---------------------------------------------------------------
% Eigennamen
%---------------------------------------------------------------
\nc{\name}[1]{\textsc{#1}\xspace}
%
\nc{\rayleigh}{\name{Rayleigh}}
\nc{\cholesky}{\name{Cholesky}}
\nc{\jacobi}{\name{Jacobi}}
\nc{\hertz}{\name{Hertz}}
\nc{\campbell}{\name{Campbell}}
\nc{\newton}{\name{Newton}}
\nc{\abbott}{\name{Abbott}}
\nc{\coulomb}{\name{Coulomb}}
\nc{\hilbert}{\name{Hilbert}}
\nc{\cattaneo}{\name{Cattaneo}}
\nc{\mindlin}{\name{Mindlin}}
\nc{\fourier}{\name{Fourier}}
\nc{\masing}{\name{Masing}}
\nc{\lms}{\name{LMS}}
\nc{\guyan}{\name{Guyan}}
\nc{\spoke}{\name{Spoke}}
\nc{\kronecker}{\name{Kronecker}}
\nc{\lanczos}{\name{Lanczos}}
\nc{\winkler}{\name{Winkler}}
\nc{\euler}{\name{Euler}}
%
%---------------------------------------------------------------
% Definition und Auswahl des Tabellenstils
%---------------------------------------------------------------
% Spaltendefinition für Tabellen (erfordert array.sty)
% Bug: zusätzliche letzte Spalte muss mit p{0mm} definiert sein (warum??)
\newcolumntype{L}{>{\raggedright}m}
\newcolumntype{R}{>{\raggedleft}m}
\newcolumntype{C}{>{\centering}m}

\nc{\dline}{\hline \hline}
\nc{\colline}[1]{\arrayrulecolor{#1} \hline}
\nc{\coldline}[1]{\arrayrulecolor{#1} \dline}

% Standardtablle mit horizontalen Linien
\nc{\tab}[5][tbh]{\begin{table}[#1]\centering\caption{#4\label{tab:#5}}\begin{tabular}{#2}\hline #3 \\ \hline\end{tabular}\end{table}}
%
% farbige Tabelle
\nc{\tabcolor}[6][tbh]{\arrayrulecolor{#6}\begin{table}[#1]\centering\caption{#4 \label {tab:#5}}\begin{tabular}{#2}\hline #3 \\\hline\end{tabular}\end{table}\arrayrulecolor{defaultcol}}
%
% Tabelle ohne Linien
\nc{\tabol}[5][tbh]{\begin{table}[#1]\centering\caption{#4\label{tab:#5}}\begin{tabular}{#2} #3 \\ \end{tabular} \end{table}}
%
% Tabelle ohne Caption und Label
\nc{\tabstern}[3][tbh]{\begin{table}[#1]\centering\begin{tabular}{#2}\hline #3 \\\hline\end{tabular}\end{table}}
%
% Buchtabelle (erfordert booktabs.sty)
\nc{\tline}{\arrayrulecolor{greyf}\toprule[1.2pt]}
\nc{\bline}{\arrayrulecolor{greyf}\bottomrule[1.2pt]}
\nc{\mline}{\arrayrulecolor{greyf}\midrule[0.8pt]}
\nc{\lline}{\arrayrulecolor{greyf}\midrule[0.5pt]}
\nc{\tabbook}[5][tb]{\begin{table}[#1]\centering\caption{#4\label{tab:#5}}\begin{tabular}{#2}\tline #3 \\ \bline\end{tabular}\end{table}}
%
% Auswahl
\rnc{\t}{\tabbook}
%
%---------------------------------------------------------------
% Figures
%---------------------------------------------------------------
%
% zentrierte Figure mit Bildunterschrift und Label
\newcommand{\f}[3][tbh]{%
 \begin{figure}[#1]
 \centering
 % \if@draft
 %    \framebox[150mm]{\raisebox{0mm}[25mm][25mm]{\texttt{#2}}}
 %\else
  \includegraphics{\chaptersubdir/#2/#2}
 %\fi
 \caption{#3}
 \label{fig:#2}
\end{figure}
}
% zentrierte Figure ohne Bildunterschrift und Label
\newcommand{\foc}[2][tbh]{%
 \begin{figure}[#1]
 \centering
  \if@draft
  \framebox[150mm]{\raisebox{0mm}[25mm][25mm]{\texttt{#2}}}
 \else
  \includegraphics{\chaptersubdir/#2/#2}
 \fi
 \label{fig:#2}
\end{figure}
}
%
% zentrierte Figure mit Bildunterschrift und Label und Skalierung
\newcommand{\fs}[4][tbh]{%
 \begin{figure}[#1]
 \centering
 \if@draft
  \framebox[150mm]{\raisebox{0mm}[25mm][25mm]{\texttt{#2}}}
 \else
  \includegraphics[scale=#4]{\chaptersubdir/#2/#2}
 \fi
 \caption{#3}
 \label{fig:#2}
\end{figure}
}
%
% zentrierte Figure in Tabelle oder auf Titelseite
% ohne Bildunterschrift, ohne Label, mit Skalierung
\newcommand{\ft}[2]{%
 \centering
 \if@draft
  \framebox[150mm]{\raisebox{0mm}[25mm][25mm]{\texttt{#1}}}
 \else
  \includegraphics[scale=#2]{\chaptersubdir/#1/#1}
 \fi
 \label{fig:#1}
}
%
% umflossene Bilder (erfordert package floatflt.sty)
%\ifnum\checkfloatflt=1
\newcommand{\ffloat}[4][p]{%
 \begin{floatingfigure}[#1]{#2}
 \centering
 \if@draft
  \framebox[150mm]{\raisebox{0mm}[25mm][25mm]{\texttt{#3}}}
 \else
  \includegraphics{\chaptersubdir/#3/#3}
 \fi
 \caption{#4}
 \label{fig:#3}
 \end{floatingfigure}
}
%\else
%\typeout{befehle.tex Warning: Package ffloat.sty not loaded!}
%\fi
%
% Zwei getrennte Figures nebeneinander
% Was ist mit Subplot?
%
\newcommand{\mf}[7][tb]{%
 \begin{figure}[#1]
 \centering
   \hs{.02\textwidth}
   \begin{minipage}{#2\textwidth}
     \centering
       \if@draft
         \framebox[150mm]{\raisebox{0mm}[25mm][25mm]{\texttt{#4}}}
       \else
         \includegraphics{\chaptersubdir/#4/#4}
       \fi
   \end{minipage}\hfill{}\hfill
   \begin{minipage}{#3\textwidth}
     \centering
       \if@draft
         \framebox[150mm]{\raisebox{0mm}[25mm][25mm]{\texttt{#5}}}
       \else
         \vs{#7}\includegraphics{\chaptersubdir/#5/#5}
       \fi
    \end{minipage}
   \hs{.02\textwidth}
   \caption{#6}
   \label{fig:#4}
 \end{figure}
}
%
%---------------------------------------------------------------
% abgesetzte Formeln
%---------------------------------------------------------------
% abgesetzte Formel mit Gleichungsnummer rechts
\nc{\e}[2]{\begin{equation} #1 \label {eq:#2} \end{equation}}
%
% lange abgesetzte Formel mit Gleichungsnummer in nächster Zeile
\nc{\el}[3][0.0ex]{\begin{samepage}
\begin{displaymath} #2 \end{displaymath} \nopagebreak[4] \vs{#1} \nopagebreak[4]\vs{-2ex} \nopagebreak[4]
\begin{equation} \label {eq:#3} \end{equation} \end{samepage}}
%
% abgesetzte Formel ohne Gleichungsnummer
\nc{\es}[2]{\begin{displaymath} #1 \label {eq:#2} \end{displaymath}}
%
% abgesetztes Formel-Array
\nc{\ea}[2]{
\begin{eqnarray}
%\setlength{\arraycolsep}{.2ex}
#1 \label {eq:#2} \end{eqnarray}}
%
% abgesetztes Formel-Array mit Gleichungsnummer in nächster Zeile
\nc{\eal}[3][0.0ex]{
\begin{samepage}
%\setlength{\arraycolsep}{.2ex}
\begin{eqnarray*}
#2
\end{eqnarray*}
\nopagebreak[4] \vs{#1} \nopagebreak[4]\vs{-2ex} \nopagebreak[4]
\begin{eqnarray}
\label {eq:#3}
\end{eqnarray}
\end{samepage}}
%
% abgesetztes Formel-Array ohne Gleichungsnummer
\nc{\eas}[2]{\begin{eqnarray*} #1 \label {eq:#2} \end{eqnarray*}}
%
%---------------------------------------------------------------
% Labels und Referenzen (deutsch/englisch)
%---------------------------------------------------------------
%
\nc{\flab}[1]{\label{fig:#1}}
\nc{\tlab}[1]{\label{tab:#1}}
\nc{\elab}[1]{\label{eq:#1}}
\nc{\slab}[1]{\label{sec:#1}}
\nc{\alab}[1]{\label{asec:#1}}
\nc{\clab}[1]{\label{ch:#1}}
\nc{\plab}[1]{\label{page:#1}}
\nc{\partlab}[1]{\label{part:#1}}
%
% Deutsch, wenn ngerman eingebunden ist (Prüfung erfordert eigenen check.sty);
% sonst englisch
%
%\ifnum\checkngerman=1
\nc{\fref}[2][]{{\bf Bild~\ref{fig:#2}{\bf #1}}}
\nc{\frefo}[2][]{{\bf \ref{fig:#2}{\bf #1}}}
\nc{\frefoo}[1]{{\bf #1}}
\nc{\frefs}[2][]{{\bf Bildern~\ref{fig:#2}{\bf #1}}}
% Tabellen
\nc{\tref}[1]{{\bf Tabelle~\ref{tab:#1}}}
\nc{\trefo}[1]{{\bf \ref{tab:#1}}}
\nc{\trefs}[1]{{\bf Tabellen~\ref{tab:#1}}}
% Gleichungen
\nc{\eref}[1]{\x{Gl.~(\ref{eq:#1})}}
\nc{\erefo}[1]{(\ref{eq:#1})}
\nc{\erefs}[1]{\x{Gln.~(\ref{eq:#1})}}
% Sections
\nc{\sref}[1]{{Abschnitt~\ref{sec:#1}}}
\nc{\srefo}[1]{\ref{sec:#1}}
\nc{\srefs}[1]{{Abschnitten~\ref{sec:#1}}}
% Sections im Anhang
\nc{\aref}[1]{{Anhang~\ref{asec:#1}}}
\nc{\arefo}[1]{\ref{asec:#1}}
\nc{\arefs}[1]{{Anhängen~\ref{asec:#1}}}
% Chapters
\nc{\cref}[1]{{Kapitel~\ref{ch:#1}}}
\nc{\crefo}[1]{\ref{ch:#1}}
\nc{\crefs}[1]{{Kapiteln~\ref{ch:#1}}}
% Pages
\nc{\pref}[1]{{Seite~\pageref{page:#1}}}
\nc{\prefo}[1]{\pageref{page:#1}}
\nc{\prefs}[1]{{Seiten~\pageref{page:#1}}}
% Parts (dies funktioniert noch nicht ?!)
\nc{\partref}[1]{{Teil~\pageref{part:#1}}}
\nc{\partrefo}[1]{\pageref{part:#1}}
\nc{\partrefs}[1]{{Teilen~\pageref{part:#1}}}
%
%---------------------------------------------------------------
% Anmerkungen am Seitenrand während des Schreibens (für fehlende Literatur usw.)
%---------------------------------------------------------------
\nc{\anm}[1]{\marginpar{\fbox{\parbox[b]{12mm}{\tiny #1}}}}
\nc{\fanm}{\marginpar{\parbox[t]{11mm}{\rule[-2.5mm]{0mm}{0mm}}} \marginpar{\fbox{\parbox[b][][c]{12mm}{\centering \scriptsize F77/F90}}}}
\nc{\fanmh}{\marginpar{\fbox{\parbox[b][][c]{11mm}{\centering \scriptsize F77/F90}}}}
\nc{\chref}{\anm{Referenz?}}
\nc{\su}{\anm{Lit. suchen!}}
\nc{\chk}{\anm{Checken!}}
\nc{\chverw}{\anm{Verweis?}}
\nc{\kom}[1]{\texttt{[#1]}}
%
%---------------------------------------------------------------
% Befehle zum Erstellen des Formelverzeichnisses
%---------------------------------------------------------------
\nc{\declare}[2]{\begin{tabular}{p{0.1\textwidth} p{0.87\textwidth}} $#1$ & #2\\ \end{tabular}}
\nc{\declarer}[2]{\begin{tabular}{p{0.1\textwidth} p{0.87\textwidth}}  \mbox{{{\rightline{ $#1$  \rule{1ex}{0ex} }}}} & #2\\ \end{tabular}}
%
%---------------------------------------------------------------
% Formeln im Text
%---------------------------------------------------------------
% Gleichung im Text; sollte statt $...$ immer verwendet werden wg. Zeilenumbruch!
\nc{\g}[1]{\x{$#1$}}
%
% Größe mit Bezeichnung nicht umbrechen
\nc{\gb}[2]{#1 \nolinebreak\g{#2}}
%
% Größe mit Bezeichnung nicht umbrechen
\nc{\bg}[2]{$#1$ \nolinebreak #2}
%
%---------------------------------------------------------------
% Größen mit Einheiten
%---------------------------------------------------------------
\nc{\geh}[2]{\x{${#1}~{\rm{[#2]}}$}}       % Größe mit Einheit nicht umbrechen
\nc{\ze}[2]{\x{${#1}~{\rm{#2}}$}}       % Größe mit Einheit nicht umbrechen
\nc{\gehg}[3]{\x{${#1}= #2 ~ {\rm{#3}}$}}       % Größe mit Einheit nicht umbrechen
\nc{\axg}[2]{\mbox{${}\si{#2}{#1}$-Achse}} % Achse mit KOS-Index
\nc{\lab}[3]{#1 \geh{#2}{#3}}
%
%---------------------------------------------------------------
% Häufige benötigte Formelelemente
%---------------------------------------------------------------
\nc{\sstar}{{+\hs{-1.78ex}\times}}      % ?
\nc{\sgn}{\mathrm{sgn}}                 % Signumfunktion
\nc{\maxf}{\mathrm{max}}                % max nicht kursiv
\nc{\minf}{\mathrm{min}}                % min nicht kursiv
\nc{\const}{\mathrm{const}}             % mathematische Konstante
\nc{\kli}{\hs{1ex}\ik}                  % ?
\nc{\itj}{{\it j}}                      % kursvives j
\nc{\iti}{{\it i}}                      % kursvives i
\nc{\itk}{{\it k}}                      % kursvives k
\nc{\itm}{{\it m}}                      % kursvives m
\nc{\itn}{{\it n}}                      % kursvives n
\nc{\bs}{\boldsymbol}                   % fette griechische Zeichen
\nc{\gr}{\ensuremath{^{\circ}}}         % Gradzeichen
\nc{\bm}{\boldmath}                     % fett
\rnc{\u}[1]{~{\mathrm{#1}}}             % steile Einheit
\nc{\intl}[2]{\int \limits_{#1}^{#2}}   % Integral mit Grenzen oben und unten
\nc{\ointl}[2]{\oint \limits_{#1}^{#2}} % Ringintegral mit Grenzen oben und unten
\nc{\suml}[2]{\sum \limits_{#1}^{#2}}   % Summe mit Indices oben und unten
\nc{\vt}{({\mathit t})}                 % von t (t)
\nc{\vtau}{(\tau)}                      % von tau (tau)
\nc{\til}{\tilde}                       % Tilde
\nc{\entspr}{~\widehat{=}~}             % 'entspricht'-Zeichen
\nc{\ub}{\unboldmath}
\nc{\bv}[2]{${}_{\rm  #1}\vec{\it e}_{\rm  #2}$}
\nc{\ds}{\displaystyle}                 % normale Schriftgröße für reduzierte Symbole
\nc{\sst}{\scriptstyle}                 % kleine Formeln
\nc{\Rt}[1]{{\mathrm{Re}} \left\{ #1 \right\}  }            % Realteil
\nc{\It}[1]{{\mathrm{Im}} \left\{ #1 \right\}  }            % Imaginärteil
\nc{\mpk}{{\scriptstyle{}_{\textstyle -}\hs{-2.1ex {{}^{}}^{(+)} }}}
\nc{\pmk}{{\scriptstyle{}_{\textstyle +}\hs{-2.1ex {{}^{}}^{(-)} }}}
\nc{\gkk}{\stackrel{(<)}{>}}            % bes. <>-Zeichen
\nc{\kgk}{\stackrel{(>)}{<}}            % bes. <>-Zeichen
\nc{\mum}{\x{\textmu m}\xspace}                  % steiles \mu für Einheiten
\nc{\hst}{\hs{-0.9ex} ^{*}}
\nc{\hhst}{\hs{-0.9ex} ^{{}^{*}}}
\nc{\hhhst}{\hs{-1.3ex} ^{{}^{*}\hs{.3ex}}}
\nc{\updots}{\mbox{\reflectbox{$\ddots$}}}  % ddots aufwärts in Formel
\nc{\diag}{\mbox{\bf diag}}             % Diagonalmatrix
\nc{\bdiag}{\mbox{\bf bdiag}}           % Block-Diagonalmatrix
\nc{\zirk}{\mbox{\bf circ}}             % zyklische Matrix
\nc{\bzirk}{\mbox{\bf bcirc}}           % Bloc-zyklische Matrix
\nc{\rz}{R\si{z}}                       % Rauheit
\nc{\ax}[2]{{}\si{{#2}}{#1}}            % Achse mit Index
\nc{\lk}{(\hs{1ex})}                    % leeres Klammerpaar
\nc{\bed}[1]{                           % geschweifte Klammer mit Bedingungen
\left\{
\begin{array}{ll}
#1
\end{array}
\right.}
%
%---------------------------------------------------------------
% spezielle häufig benutzte Größen
%---------------------------------------------------------------
%\nc{\mhq}[4]{\m{\hat{\Phi \hs{-.3ex}}}{#1}{#2}{ #3}{#4}}
\nc{\fc}{F\si{C}}
\nc{\eip}{e^{\ii \, \Delta \phi}}
\nc{\emip}{e^{-\ii\, \Delta \phi}}
\nc{\om}[2]{\ubm{{}_{\hspace{.6ex} \rm #1}^{\rm #2}}{\omega}}
\nc{\dotom}[2]{\ubm{{}_{\hspace{.6ex} \rm #1}^{\rm#2}}\dot{\omega}}
\nc{\tildeom}[2]{\ubm{{}_{\hspace{.6ex} \rm #1}^{\rm #2}}\tilde{\omega}}
\nc{\eo}{{\mathit{EO}}}
\nc{\nd}{{\mathit{ND}}}
\nc{\ome}{\Omega\si{E}}
\nc{\omr}{\Omega\si{R}}
\nc{\omp}{\Omega\si{P}}
\nc{\omn}{\Omega\si{0}}
\nc{\dphi}{\Delta \phi}
\nc{\ns}{{\it n}\si{S}}
\nc{\nkonp}{{\it n}\si{K, FEM}}
\nc{\nh}{{\it n}\si{H}}
\nc{\ndis}{{\it n}\si{Dis}}
\nc{\nms}{{\it n}\si{M,S}}
\nc{\nmd}{{\it n}\si{M,D}}
\nc{\nfem}{{\it n}\si{FEM}}
\nc{\nfhg}{{\it n}\si{FHG}}
\nc{\nl}{{\it n}\si{L}}
\nc{\nseg}{{\it n}\si{Seg}}
\nc{\nm}{{\it n}\si{M}}
\nc{\cnn}{c\si{N0}}
\nc{\ctn}{c\si{T0}}
\nc{\mopt}{m\si{Opt}}
\nc{\ures}{\hat{u}\si{{\it x}, Res}}
\nc{\fres}{f\si{E, Res}}
\nc{\tp}{M\si{r}}
\nc{\deltah}{\frac{\delta}{2}}
\nc{\fur}{\mathrm{f"ur~}}
\nc{\OL}{{\rm OL}}
\nc{\OR}{{\rm OR}}
\nc{\OLp}{{\rm OL}' }
\nc{\ORp}{{\rm OR}' }
\nc{\matdim}[3]{\g{[#1\times #2]}-#3}
\nc{\matdims}[2]{\g{[#1\times #2]}}
%
%---------------------------------------------------------------
% Abstände vor Satzzeichen nach und in Formeln
%---------------------------------------------------------------
\nc{\fabstand}{\,}
\nc{\fp}{\fabstand.}
\nc{\fk}{\fabstand,}
\nc{\fsk}{\fabstand;}
\nc{\fdp}{\fabstand:}
%
%---------------------------------------------------------------
% spezielle mathematische Symbole (kalligraphisches M, D, K, H, S, F)
%---------------------------------------------------------------
\nc{\dsym}{\mm{\mathcal D}}
\nc{\fsym}{\mm{\mathcal F}}
\nc{\hsym}{\mm{\mathcal H}}
\nc{\ksym}{\mm{\mathcal K}}
\nc{\msym}{\mm{\mathcal M}}
\nc{\ssym}{\mm{\mathcal S}}
\nc{\hmsym}{\hat{\msym}}
\nc{\hksym}{\hat{\ksym}}
\nc{\hdsym}{\hat{\dsym}}
\nc{\hhsym}{\hat{\hsym}}
%
%---------------------------------------------------------------
% angepasste Klammern
%---------------------------------------------------------------
\nc{\la}{\left(}
\nc{\ra}{\right)}
\nc{\lb}{\left[}
\nc{\rb}{\right]}
\nc{\lc}{\left\{}
\nc{\rc}{\right\}}
\nc{\ld}{\left.}
\nc{\rd}{\right.}
\nc{\vr}{\rule{0mm}{2ex}}
%
%---------------------------------------------------------------
% Ersetzen von phi, theta, rho und epsilon
%---------------------------------------------------------------
\rnc{\phi}{\varphi}
\rnc{\theta}{\vartheta}
\rnc{\rho}{\varrho}
\rnc{\epsilon}{\varepsilon}
%
%---------------------------------------------------------------
% steile mathematische Symbole
%---------------------------------------------------------------
\nc{\dd}{{\mathrm{d}}}
\nc{\ii}{{\mathrm{i}}}
\nc{\ee}{{\mathrm{e}}}
%
%---------------------------------------------------------------
% inverse, transponierte und allgemeine Indices
%---------------------------------------------------------------
\nc{\inv}{^{-1}}                    % invers
\nc{\tra}{^{\mathrm T}}                 % transponiert
\nc{\herm}{^{\mathrm H}}                % hermitesch transponiert
\nc{\stern}{^{*}}                      % Stern
\nc{\stst}{^{**}}                   % Sternstern
\nc{\si}[1]{_{\mathrm{#1}}}              % steiler tiefgestellter Index
\nc{\ki}[1]{_{\mathit{#1}}}              % kursiver tiefgestellter Index
%
%---------------------------------------------------------------
% kursiver hochgestellter Zählindex (i)
%---------------------------------------------------------------
\nc{\ik}{({\it i})}
\nc{\jk}{({\it j})}
\nc{\kk}{({\it k})}
\nc{\nk}{{\it n}}
\nc{\ikp}{({\it i}+1)}
\nc{\ikm}{({\it i}-1)}
\nc{\jkp}{({\it j}+1)}
\nc{\jkm}{({\it j}-1)}
\nc{\kkp}{({\it k}+1)}
\nc{\kkm}{({\it k}-1)}
\nc{\nkp}{({\it n}+1)}
\nc{\nkm}{({\it n}-1)}
\nc{\mkp}{({\it m}+1)}
\nc{\mkm}{({\it m}-1)}
%
%---------------------------------------------------------------
% Matrizen und Vektoren
%---------------------------------------------------------------
\rnc{\matrix}[2]{\left[\!\!\begin{array}{#1} #2\end{array}\!\!\right]}
\rnc{\vector}[1]{\matrix{c}{#1}}
%---------------------------------------------------------------
%
% Matrix und Vektor fett kursiv
\nc{\mm}[1]{\boldsymbol{#1}}
% Matrix und Vektor fett steil
%\nc{\mm}[1]{\mathbf{#1}}
%
% Vektoren und Matrizen ohne Index fett kursiv
\nc{\m}[5]{{}_{\mathrm{#2}}^{\mathrm{#3}}{\mm{#1}}^{\mathrm{#4}}_{\mathrm{#5}}}
%
% Vektoren und Matrizen mit Index fett steil
%\nc{\m}[5]{{}_{\mathrm{#2}}^{\mathrm{#3}}{\mathbf{#1}}^{\mathrm{#4}}_{\mathrm{#5}}}
%
%
%\nc{\mm}[1]{\boldmath{#1}} % Vektoren und Matrizen ohne Index fett kursiv
%
%\nc{\mm}[1]{\mathbf{#1}} % Vektoren und Matrizen ohne Index fett steil
%
% Besondere Matrizen:
\nc{\mnull}{\mm{0}}
\nc{\mma}{\mm{A}}
\nc{\mmb}{\mm{B}}
\nc{\mmc}{\mm{C}}
\nc{\mmd}{\mm{D}}
\nc{\mme}{\mm{E}}
\nc{\mmf}{\mm{F}}
\nc{\mmg}{\mm{G}}
\nc{\mmh}{\mm{H}}
\nc{\mmi}{\mm{I}}
\nc{\mmj}{\mm{J}}
\nc{\mmk}{\mm{K}}
\nc{\mml}{\mm{L}}
\nc{\mmm}{\mm{M}}
\nc{\mmn}{\mm{N}}
\nc{\mmo}{\mm{O}}
\nc{\mmp}{\mm{P}}
\nc{\mmq}{\mm{Q}}
\nc{\mmr}{\mm{R}}
\nc{\mms}{\mm{S}}
\nc{\mmt}{\mm{T}}
\nc{\mmu}{\mm{U}}
\nc{\mmv}{\mm{V}}
\nc{\mmw}{\mm{W}}
\nc{\mmx}{\mm{X}}
\nc{\mmy}{\mm{Y}}
\nc{\mmz}{\mm{Z}}
%
\nc{\mphi}{\mm{\phi}}
\nc{\mtildephi}{\mm{\tilde{\phi}}}
\nc{\mpsi}{\mm{\psi}}
\nc{\mhatphi}{\mm{\hat \phi}}
\nc{\mhatpsi}{\mm{\hat \psi}}
\nc{\mtildepsi}{\mm{\tilde \psi}}
\nc{\hq}{\hat q}
\nc{\mhq}{\mm{\hat q}}
%
\nc{\skao}[1]{#1}           % Definition eines Skalars ohne Indices
\nc{\skai}[5]{{}_{\rm  #2}^{\rm #3}{#1}^{\rm #4}_{\rm  #5}}   % Definition eines Skalars mit Indices
\nc{\veco}[1]{\vec{#1}}     % Definition eines Vektors mit Vektorpfeil ohne Indices
\nc{\veci}[5]{{}_{\rm  #2}^{\rm #3}{\vec{#1}}^{\rm #4}_{\rm  #5}} % Definition eines Vektors mit Vektorpfeil mit Indices
\nc{\vectupo}[1]{\vec{\bf{#1}}}
\nc{\vectupi}[5]{{}_{\rm  #2}^{\rm #3}{\vec{\bf{#1}}}^{\rm #4}_{\rm#5}}
%
%---------------------------------------------------------------
% trigon. Funktionen
%---------------------------------------------------------------
\nc{\calpha}{\cos \alpha}
\nc{\salpha}{\sin \alpha}
\nc{\talpha}{\tan \alpha}
\nc{\cbeta}{\cos \beta}
\nc{\sbeta}{\sin \beta}
\nc{\tbeta}{\tan \beta}
\nc{\cgamma}{\cos \gamma}
\nc{\sgamma}{\sin \gamma}
\nc{\tgamma}{\tan \gamma}
\nc{\cdelta}{\cos \delta}
\nc{\sdelta}{\sin \delta}
\nc{\tdelta}{\tan \delta}
\nc{\ctheta}{\cos \theta}
\nc{\stheta}{\sin \theta}
\nc{\ttheta}{\tan \theta}
\nc{\sphi}{\sin \phi}
\nc{\cphi}{\cos \phi}
\nc{\tphi}{\tan \phi}
\nc{\spsi}{\sin \psi}
\nc{\cpsi}{\cos \psi}
\nc{\tpsi}{\tan \psi}
%
%---------------------------------------------------------------
% Blankzeichen
%---------------------------------------------------------------
%\nc{\bl}{\hs{-1.4ex}\raisebox{-0.9ex}{\resizebox{1.3ex}{1ex}{\rotatebox{90}{[}}}\hs{-1.3ex}}
\nc{\bl}{\textvisiblespace} % funktioniert nur ohne T1 encoding richtig!
%
%---------------------------------------------------------------
% Ersetzungen durch psfrag
% Sollte nicht mehr im Dokument verwendet werden, sondern schon bei
% der Erzeugung des Bildes!
%---------------------------------------------------------------
%\nc{\p}[1]{\psfrag{#1}[c][c]{\figfont #1}}
%\nc{\pf}[2]{\psfrag{#1}[c][c]{\figfont #2}}
%
%---------------------------------------------------------------
% spezielle 'Verschönerungen', schriftabhängig
%---------------------------------------------------------------
%
\nc{\hfsym}{\hat{\fsym\hs{0.5ex}}\hs{-.5ex}}
\nc{\hssym}{\hat{\ssym\hs{0.5ex}}\hs{-.5ex}}
\nc{\mPhi}[4]{\m{\Phi}{#1}{#2}{\hs{-.4ex} #3}{\hs{-.5ex} #4}}
\nc{\mDeltaPhi}[4]{\m{\Delta \Phi}{#1}{#2}{\hs{-.4ex} #3}{\hs{-.5ex} #4}}
\nc{\mPsi}[4]{\m{\Psi}{#1}{#2}{\hs{0ex} #3}{\hs{0ex} #4}}
\nc{\mtildePhi}[4]{\m{\tilde{\Phi \hs{-.3ex}}}{#1}{#2}{ #3}{#4}}
\nc{\mhatPhi}[4]{\m{\hat{\Phi \hs{-.3ex}}}{#1}{#2}{ #3}{#4}}
\nc{\mtildePsi}[4]{\m{\tilde{\Psi \hs{0ex}}}{#1}{#2}{ #3}{#4}}
\nc{\mhatPsi}[4]{\m{\hat{\Psi \hs{.2ex}}}{\hs{-.2ex} #1}{\hs{-.2ex} #2}{\hs{-.2ex} #3}{\hs{-.2ex} #4}}
%
%---------------------------------------------------------------
% Sections usw. automatisch mit label versehen
%---------------------------------------------------------------
\nc{\partl}[2]{\part{#1 \partlab{#2}}}
\nc{\chapterl}[2]{\chapter{#1 \clab{#2}}}
\nc{\sectionl}[2]{\section{#1 \slab{#2}}}
\nc{\subsectionl}[2]{\subsection{#1 \slab{#2}}}
\nc{\subsubsectionl}[2]{\subsubsection{#1 \slab{#2}}}
%
\nc{\partls}[3]{\part[#1]{#2 \partlab{#3}}}
\nc{\chapterls}[3]{\chapter[#1]{#2 \clab{#3}}}
\nc{\sectionls}[3]{\section[#1]{#2 \slab{#3}}}
\nc{\subsectionls}[3]{\subsection[#1]{#2 \slab{#3}}}
\nc{\subsubsectionls}[3]{\subsubsection[#1]{#2 \slab{#3}}}
%
\nc{\ch}[2]{\chapter{#1 \clab{#2}} \thispagestyle{kapitelbeginn}}
\nc{\cha}[2]{\chaptera{#1 \alab{#2}}}
%\nc{\cha}[2]{\chapter{#1 \alab{#2}} \markright{\headertextfont Anhang}}
%\nc{\cha}[2]{\chapter{#1 \alab{#2}} \markright{\headertextfont Anhang} \thispagestyle{kapitelbeginn}}
\nc{\sect}[2]{\section{#1 \slab{#2}}}
\nc{\ssect}[2]{\subsection{#1 \slab{#2}}}
\nc{\sssect}[2]{\subsubsection{#1 \slab{#2}}}
\nc{\chs}[2]{\chapter*{#1 \clab{#2}} \thispagestyle{kapitelbeginn}}
%
%---------------------------------------------------------------
% Aufzählungen usw.
%---------------------------------------------------------------
\nc{\bitemize}{\begin{itemize}}
\nc{\eitemize}{\end{itemize}}
\rnc{\labelitemi}{$\bullet$}
\rnc{\labelitemii}{$\circ$}
%
\nc{\benumerate}{\begin{enumerate}}
\nc{\eenumerate}{\end{enumerate}}
\nc{\barray}{\begin{array}}
\nc{\earray}{\end{array}}
\nc{\bsamepage}{\begin{samepage}}
\nc{\esamepage}{\end{samepage}}
%---------------------------------------------------------------
% Weitere Abkürzungen
%---------------------------------------------------------------
\nc{\ols}{$\x{OL}'$}
\nc{\ors}{$\x{OR}'$}
%
%---------------------------------------------------------------
% Fußnoten
%---------------------------------------------------------------
%
\nc{\fnote}[1]{\footnote{\samepage #1}} % Fußnoten nicht umbrechen!
\nc{\fnotetext}[1]{\footnotetext{\samepage #1}} % Fußnoten nicht umbrechen!
%(Dies führt u.U. zu seltsamen Formatierungen, checken! -> umgebrochene Fußnoten sehen aber richtig schlimm aus!)
%
%---------------------------------------------------------------
% Indexerstellung
%---------------------------------------------------------------
% Übernehmen von Indexeinträgen aus dem Text - Originalbefehl \index{}
\nc{\ind}[1]{#1\index{#1}}                  % Index direkt aus Text anlegen
\nc{\indsub}[3][]{#2\index{#3!#2#1}}        % Subindex im Text anlegen (mit optionaler Endung)
\nc{\indsubsub}[4][]{#2\index{#3!#4!#2#1}}  % Subsubindex im Text anlegen (mit optionaler Endung)
%
%---------------------------------------------------------------
% Farbdefinition
% Erfordert package color.sty oder colortbl.sty
%---------------------------------------------------------------
%\ifnum\checkcolor=1
\definecolor{red}{rgb}{1.0,0.0,0.0}
\definecolor{green}{rgb}{0.0,1.0,0.0}
\definecolor{blue}{rgb}{0.0,0.0,1.0}
\definecolor{black}{rgb}{0.0,0.0,0.0}
\definecolor{white}{rgb}{1.0,1.0,1.0}
\definecolor{yellow}{rgb}{1.0,1.0,0.0}
\definecolor{orange}{rgb}{1.0,0.4,0.0}
\definecolor{purple}{rgb}{0.6,0.3,0.5}
\definecolor{cyan}{rgb}{0.0,1.0,1.0}
\definecolor{magenta}{rgb}{1.0,0.0,1.0}
\definecolor{greyi}{rgb}{0.9,0.9,0.9}
\definecolor{greyh}{rgb}{0.8,0.8,0.8}
\definecolor{greyg}{rgb}{0.7,0.7,0.7}
\definecolor{greyf}{rgb}{0.6,0.6,0.6}
\definecolor{greye}{rgb}{0.5,0.5,0.5}
\definecolor{greyd}{rgb}{0.4,0.4,0.4}
\definecolor{greyc}{rgb}{0.3,0.3,0.3}
\definecolor{greyb}{rgb}{0.2,0.2,0.2}
\definecolor{greya}{rgb}{0.1,0.1,0.1}
\definecolor{defaultcol}{rgb}{0.0,0.0,0.0}
\nc{\red}[1]{\textcolor{red}{#1}}
\nc{\blue}[1]{\textcolor{blue}{#1}}
\nc{\green}[1]{\textcolor{green}{#1}}
\nc{\yellow}[1]{\textcolor{yellow}{#1}}
\nc{\black}[1]{\textcolor{black}{#1}}
\nc{\white}[1]{\textcolor{white}{#1}}
\nc{\violet}[1]{\textcolor{purple}{#1}}
\nc{\orange}[1]{\textcolor{orange}{#1}}
\nc{\cyan}[1]{\textcolor{cyan}{#1}}
\nc{\magenta}[1]{\textcolor{magenta}{#1}}
\nc{\greya}[1]{\textcolor{greya}{#1}}%
\nc{\greyb}[1]{\textcolor{greyb}{#1}}
\nc{\greyc}[1]{\textcolor{greyc}{#1}}
\nc{\greyd}[1]{\textcolor{greyd}{#1}}
\nc{\greye}[1]{\textcolor{greye}{#1}}
\nc{\greyf}[1]{\textcolor{greyf}{#1}}
\nc{\greyg}[1]{\textcolor{greyg}{#1}}
\nc{\greyh}[1]{\textcolor{greyh}{#1}}
\nc{\greyi}[1]{\textcolor{greyi}{#1}}
%\else
%\typeout{befehle.tex Warning: Package color.sty not loaded!}
%\fi
%------------------------------------------------------------
\nc{\musteil}{\textmu\xspace}
